\scanpage{15}
\uline{Cohomologies relatives} (un cas spécial)

1. \uline{Caractérisation des foncteurs} $A\longrightarrow\mathrm{Hom}(A,B)$.

Soit $C$ une catégorie. Pour tout $B\in C$, soit $h_B$ le foncteur de $C$ dans les ensembles tel que
\[
 h_B(A)=\mathrm{Hom}(A,B)
\]
$h_B$ est un foncteur \unsure{.l|,!|,} contravariant en $B$, à valeurs dans la catégorie des foncteurs $C\longrightarrow\unsure{\mathcal E|E}$, et on a
\[
 \alpha:\mathrm{Hom}_C(B,B')\xrightarrow{\sim}\mathrm{Hom}(h_B,h_{B'}) \tag*{(*)}
\]
[En effet, on définit la flèche inverse $\beta$, en associant à \deleted{\unsure{tout}} $u\in\mathrm{Hom}(h_B,h_{B'})$ \deleted{\unsure{un}} l'élément $\beta(u)$ de $\mathrm{Hom}_C(B,B')$ image par $u_B$ de l'élément unité de $h_B(B)=\mathrm{Hom}_{\unsure{C|a}}(B,B)$ dans $h_B(B')=\mathrm{Hom}_C(B,B')$].

\deleted{Si dans $C$ il y a \unsure{assez de sommes}} Si $C$ est additive, alors $h_B$ est additif \inserted{\uline{exact à g.}}, le foncteur $B\longrightarrow h_B$ est additif (exact \uline{à gauche}) et l'isomorphisme $(*)$ est valable en se limitant aux hom. \uline{additifs}. Ainsi $B\longrightarrow h_B$ \deleted{\unsure{est}} permet d'``identifier'' $C$ à une sous catégorie de la catégorie $\mathrm{Fonc.add.}(C^o,\unsure{\mathscr G|\mathfrak G|\mathcal Y})$.

Supposons que $C$ ait un générateur $U$, et \deleted{que} satisfasse à AB 3, 4, 5, . \deleted{\unsure{Comme|Le foncteur}} \inserted{\unsure{Tout|Tjs est}} \deleted{\unsure{alors}} $h_B$ transforme limites inductives en limites projectives, et \deleted{\unsure{par suite} \unsure{aussi}} \inserted{est exact} à gauche. On va donner une réciproque. Tout d'abord, si $h,h'$ sont deux foncteurs \inserted{contr.} $C^{\deleted{\unsure{o}}}\longrightarrow\mathscr G$ ayant
% archive sheet number: 14
